\section*{Chapter \ref{ch:iv2:making-a-market} --- Making a Market}
\markboth{\protect\omTocMark{}{Answers: Making a Market}}{\protect\omTocMark{}{Answers: Making a Market}}

\begin{solution}{iq:iv2:making-a-market:1}
Quote 10 bid, 11 asked, and expect to earn 0.5. A quote of 11 bid, 12 asked also earns 0.5 on average.

Each die averages 3.5, so the sum is worth 10.5. If the customer sells to you at 10, you gain 0.5 on average, and if they buy at 11, you gain 0.5. At 11 and 12, buying at 11 loses 0.5 and selling at 12 gains 1.5, and the average of the two is again 0.5. A customer whose side is a coin toss pays you half your width wherever the quote sits.

\iqwrong{Concluding that the position of a quote does not matter. It does not matter only because this customer ignores it. Someone who has worked out 10.5 would sell to you at 11 every time, and you would lose 0.5 on each trade.}
\iqfollow{How much risk are you taking for that 0.5?}{The sum has a standard deviation of $\sqrt{35/4} \approx 2.96$, about six times your edge. One trade proves nothing. The edge shows up only over many trades.}
\iqwhy{Width pays you for uninformed flow, and the centre protects you against everyone else. Say both when you give a quote.}
\end{solution}

\begin{solution}{iq:iv2:making-a-market:2}
Quote 6 bid, 8 asked, and lose $30/13 \approx 2.31$ on average. The narrowest whole-number width that loses less than 1 is 6, at 4 bid, 10 asked, which loses $12/13 \approx 0.92$.

The ranks are equally likely and centred on 7, so the symmetric quote is best. At 6 and 8, the ranks 1 to 5 sell to you and lose you 5, 4, 3, 2 and 1, and the ranks 9 to 13 buy from you and lose you 1 to 5. That totals 30 over 13 ranks. The best quote at each width loses:
\begin{center}\small
\begin{tabular}{@{}lccccc@{}}
Width & 2 & 4 & 5 & 6 & 12 \\ \hline
Average loss & 2.31 & 1.54 & 1.23 & 0.92 & 0
\end{tabular}
\end{center}

\iqwrong{Looking for a quote that makes money. Someone who knows the rank trades only when you are wrong, so you can never profit. Only at 1 bid, 13 asked do you stop losing, and there nobody trades. Width is a price you set on being picked off.}
\iqwhy{This is the extreme case of adverse selection\index{adverse selection}. When you suspect the other side knows, the question stops being ``what is it worth?'' and becomes ``how much am I willing to pay to find out?''}
\end{solution}

\begin{solution}{iq:iv2:making-a-market:3}
Quote 14 bid, 16 asked. The contract is worth $91/6 \approx 15.17$, not $3.5^2 = 12.25$.

The expected value is $(1 + 4 + 9 + 16 + 25 + 36)/6 = 91/6$. Of the half-steps, a mid of 15 is the nearest to 15.17 that fits a width of 2, so the quote is 14 to 16. The customer's side is a coin toss, so the edge is 1 on average.

\iqwrong{Squaring the average die. The expectation of a square is the square of the expectation plus the variance, and the die's variance is $35/12 \approx 2.92$, which is exactly the gap between 12.25 and 15.17. This is Jensen's inequality\index{Jensen's inequality}: a convex payoff is worth more than the payoff at the average.}
\iqfollow{What about a contract paying 1 divided by the die?}{It is also convex in the die, so again the average payoff beats the payoff at the average: $49/120 \approx 0.41$ against $1/3.5 \approx 0.29$.}
\end{solution}

\begin{solution}{iq:iv2:making-a-market:4}
Quote the next card 5.75 bid, 6.75 asked, and the sum of the rest 263.5 bid, 264.5 asked. The first carries real risk, and the second carries none.

Each rank from 1 to 13 appears four times, so the whole deck totals $4 \times 91 = 364$. With 100 showing, the 42 unseen cards total exactly 264. Each is equally likely to be next, so the next card is worth $264/42 = 44/7 \approx 6.29$, and the nearest quarter gives a mid of 6.25. Your counterparty's value lies inside both quotes, so they pass.

\iqwrong{Treating the sum as uncertain because it adds 42 random cards. Those cards are not independent. The total is fixed, so the sum's variance is zero. Averaging $42 \times 7 = 294$ ignores what you have seen.}
\iqfollow{What would you do if someone offered to sell you the sum at 263?}{Buy as much as they will sell, because you make 1 on each lot for certain. Anyone who can count the deck will not offer it.}
\end{solution}

\begin{solution}{iq:iv2:making-a-market:5}
They sell you both $A$ at 0.55 and $B$ at 0.50, collecting 1.05 and owing exactly 1, for a sure profit of 0.05. To rule this out, the two bids must add up to at most 1, and the two asks must add up to at least 1.

Exactly one of $A$ and $B$ pays, because the two events cover every result of the die with no overlap, so a pair of one lot each is worth exactly 1 whatever happens. A pair sold to you for more than 1 is free money for the other side. A pair bought from you for less than 1 is free money too. Your asks here add up to 1.12, so that side is safe.

\iqfollow{Is 0.40 to 0.45 on both safe?}{No. The asks total 0.90, so they buy both from you for 0.90 and receive 1, making 0.10 for sure.}
\end{solution}

\begin{solution}{iq:iv2:making-a-market:6}
Quote 4 bid, 5 asked before the rolls, and 3.5 bid, 4.5 asked after the 3. After a 6, the contract is worth exactly 6, so you would show 5.5 bid, 6.5 asked.

The higher of two dice equals $k$ in $2k - 1$ of the 36 outcomes, so the value is $\sum k(2k - 1)/36 = 161/36 \approx 4.47$. Given a first die of 3, the second die gives a maximum of 3 in three cases, and otherwise 4, 5 or 6, so the value is $(3 + 3 + 3 + 4 + 5 + 6)/6 = 4$.

After the first die the values are:
\begin{center}\small
\begin{tabular}{@{}l*{6}{c}@{}}
First die & 1 & 2 & 3 & 4 & 5 & 6 \\ \hline
Value & 3.50 & 3.67 & 4.00 & 4.50 & 5.17 & 6.00
\end{tabular}
\end{center}

\iqfollow{Why would you ever keep a width after the 6?}{You need none for risk, because nothing is uncertain. The stem makes you show one, and the customer's coin toss pays you for it. A real customer would not trade at 6.5 for something worth 6.}
\end{solution}

\begin{solution}{iq:iv2:making-a-market:7}
Quote about 159 bid, 179 asked. There are 168 primes below 1,000, so that quote passes.

The prime number theorem\index{prime number theorem} says the count of primes up to $n$ is about $n/\ln n$, and a sharper form is about $n/(\ln n - 1)$. With $\ln 1000 = 3 \times 2.303 \approx 6.91$, the sharper form gives $1000/5.91 \approx 169$. The centre of the quote is 169, and 20 wide makes it 159 to 179.

If you do not recall the theorem, sieve. About $1000 \times \tfrac12 \times \tfrac23 \times \tfrac45 \times \tfrac67 \approx 229$ numbers have no factor of 2, 3, 5 or 7, and the exact count is 228. Those still include composites such as $11 \times 11$, so the answer is well below 229.

\iqwrong{Using $n/\ln n \approx 145$ as it stands. The quote 135 to 155 then loses $168 - 155 = 13$ when the interviewer buys at 155. The plain form runs low for small $n$.}
\end{solution}

\begin{solution}{iq:iv2:making-a-market:8}
Quote 5 bid, 7 asked. Both sides earn 1 on average, and they differ completely in shape. If you are long at 5, you lose money with probability $1 - (5/6)^4 \approx 52\%$, and at most 4. If you are short at 7, you lose with probability $(5/6)^7 \approx 28\%$, and the loss has no limit.

The number of rolls is geometric\index{geometric distribution} with mean 6 and variance 30, so the standard deviation is about 5.5. Its median is 4, since $1 - (5/6)^3$ is below one half and $1 - (5/6)^4$ is above. Most results are below the mean, and a few are far above it.

So a long position usually loses a little and sometimes wins a lot, and a short position usually wins a little and sometimes loses a lot. It takes 20 rolls or more about 3\% of the time, and then a short loses at least 13.

\iqwrong{Centring on the median because ``most outcomes are below 6''. The contract pays the number of rolls, so the mean is what breaks even over many trades.}
\iqfollow{Which side would you rather have at the same average profit?}{It depends on your limits. A short has a small chance of a very large loss, so it needs a limit, more edge, or both. Say that before you say which you prefer.}
\iqwhy{Selling a skewed payoff at a fair average is how selling options looks. A position may have a good expected value and still be one you cannot afford to hold.}
\end{solution}

\begin{solution}{iq:iv2:making-a-market:9}
Quote 4 bid, 6 asked. The result lands outside the quote with probability about 29\%. After a first red card, the value is $92/17 \approx 5.41$, and you quote 4.5 bid, 6.5 asked.

By symmetry each card is red with probability one half, so the expected number of reds is 5 whatever the dependence between cards. The count is hypergeometric\index{hypergeometric distribution}, with variance
\[
10 \cdot \tfrac12 \cdot \tfrac12 \cdot \tfrac{52 - 10}{52 - 1} = \tfrac{35}{17},
\]
so the standard deviation is 1.43, below the 1.58 of ten independent coin tosses. Drawing without replacement pulls the count towards the middle. Summing the exact probabilities of 4, 5 and 6 gives 0.71.

After a red, the other nine cards come from 25 red in 51, so the value is $1 + 9 \cdot 25/51$. The first card was worth 0.5 and is now worth 1, a gain of 0.5. The other nine fall from 4.5 to 4.41, because a red has left the deck. Together the value rises by 0.41.

\iqwrong{Moving up by 1, or by 0.5. A rise of 1 forgets that the first card already counted for 0.5. A rise of 0.5 forgets that the other nine cards depend on the first.}
\end{solution}

\begin{solution}{iq:iv2:making-a-market:10}
With $k = h/\sigma$, the expected loss is $2\sigma\left[\varphi(k) - k\,\Phi(-k)\right]$, where $\varphi$ and $\Phi$ are the standard normal density and distribution function. A 90\% interval, $k = 1.645$, loses $0.042\,\sigma$. To lose less than $0.05\,\sigma$, you need $k \ge 1.57$, a width of about $3.1\,\sigma$.

The counterparty buys when $X > \mu + h$, costing you $X - \mu - h$, and sells in the mirror case. So the loss is
\[
2\,\mathbb{E}\left[(X - \mu - h)^+\right] = 2\sigma\int_k^\infty (z - k)\,\varphi(z)\,dz = 2\sigma\left[\varphi(k) - k\,\Phi(-k)\right].
\]
This is the normal loss function\index{normal loss function}. At $k = 0$ it is $0.80\,\sigma$, and at $k = 1$ it is $0.167\,\sigma$.

\iqwrong{Thinking a 90\% interval loses 10\% of something. You trade 10\% of the time, but those trades lose only a small amount on average, because the tail is thin just beyond the interval.}
\iqfollow{Why not quote very wide, then?}{Against this counterparty you should. But you also want the uninformed flow, and it stops trading if you are wide. How wide you go is a balance between those two, not just a statement of your uncertainty.}
\end{solution}

\begin{solution}{iq:iv2:making-a-market:11}
Centre on about 254. The contract pays the number, so the fair value is the mean, and for a lognormal\index{lognormal distribution} the mean is the median times $e^{\sigma^2/2}$. Here $\sigma = \ln 2$, so the factor is $e^{0.24} \approx 1.27$.

So $200 \times 1.27 \approx 254$. The true number exceeds 254 with probability $\Phi(-\sigma/2) = \Phi(-0.35) \approx 36\%$. The 90\% range is $200 \cdot 2^{\pm 1.645}$, about 64 to 625.

\iqwrong{Centring on the best guess of 200. Half the time the number is above 200, but it can be far above, while it can only be a little below in absolute terms. If you sell at 200 again and again, you lose 54 a lot on average. A quote centred on the median gets the shape wrong as well: the range is not symmetric around the centre.}
\iqfollow{How did you get a factor of 2?}{A Fermi estimate is a product of factors, and logs of a product add. Three factors each unsure by a factor of 1.5 give a log standard deviation of $\sqrt3 \ln 1.5$, which is a factor of about 2.0 in all.}
\iqwhy{Unknown quantities are usually products, so their uncertainty is multiplicative. Your bid can be far nearer zero than your ask is to your centre, and saying so is a mark of someone who has thought about it.}
\end{solution}

\begin{solution}{iq:iv2:making-a-market:12}
Quote 94 bid, 106 asked. The break-even ask is the expected value given that someone buys, and the break-even bid is the expected value given that someone sells.

Use Bayes\index{Bayes' rule}. If the value is 120, a buy comes with probability $0.30 + 0.70 \cdot \tfrac12 = 0.65$. If it is 80, a buy comes only from the coin toss, $0.70 \cdot \tfrac12 = 0.35$. So given a buy, the chance of 120 is $0.65/(0.65 + 0.35) = 0.65$, and the ask is $80 + 40 \times 0.65 = 106$. The bid is the mirror image, 94.

So the width is the share of knowing traders times the range, $0.30 \times 40 = 12$. This is the Glosten-Milgrom\index{Glosten-Milgrom model} quote. Uninformed traders pay exactly what the informed ones take from you.

\iqwrong{Quoting around 100 with a width picked by feel. What matters is the value given that someone buys. There is a quick route here: 30\% of buyers know and pay for 120, and the rest pay for an average of 100, so the ask is $0.3 \times 120 + 0.7 \times 100 = 106$. It works only because the prior is one half. With any other prior, the share of knowers among buyers differs from 30\%.}
\iqfollow{What if everyone knows, or nobody does?}{With everyone knowing, the quote is 80 and 120, and the trader gains nothing at those prices. With nobody knowing, it is 100 and 100, with no width at all, as there is nothing to protect against in this model.}
\end{solution}

\begin{solution}{iq:iv2:making-a-market:13}
The quotes before the five traders are below, and the order does not matter, because each purchase multiplies the odds of 100 by 1.5, and each sale divides them by 1.5.
\begin{center}\small
\begin{tabular}{@{}lccccc@{}}
Trader & 1 & 2 & 3 & 4 & 5 \\ \hline
Trades so far & none & B & BB & BBS & BBSB \\
Bid & 40.00 & 50.00 & 60.00 & 50.00 & 60.00 \\
Ask & 60.00 & 69.23 & 77.14 & 69.23 & 77.14
\end{tabular}
\end{center}

If the value is 100, a buy has probability $0.20 + 0.40 = 0.60$, and if it is 0, it has probability $0.40$. The likelihood ratio\index{likelihood ratio} is 1.5, and a sale gives its reciprocal. So with $n$ net purchases the odds of 100 are $1.5^n$. The ask is the value given one more purchase, at odds $1.5^{n+1}$, and the bid is the value given one more sale, at odds $1.5^{n-1}$. With $n = 1$ the ask is $2.25/3.25$ of 100, which is 69.23. The final odds are $1.5^2$, so the chance of 100 is $9/13$.

\iqfollow{What if uninformed traders buy 70\% of the time?}{A buy now has probability $0.76$ or $0.56$, a ratio of 1.36, and a sale has ratio $0.24/0.44 = 0.55$. Since $1/1.36 = 0.74$, a sale moves you more than a buy: the rarer trade says more. Order still does not matter, but the net count is no longer enough, and you need both counts.}
\end{solution}

\begin{solution}{iq:iv2:making-a-market:14}
The contract is worth 2.7 after the sale, and you quote 2.2 bid, 3.2 asked. After a pass, it is worth 3.5.

Weigh each value by the probability of a sale. For values 1 and 2 the looker sells, so that is $\tfrac12 + \tfrac12 \cdot \tfrac12 = \tfrac34$. For values 3 to 6 only the coin toss can sell, so $\tfrac14$. The weights sum to $\tfrac52$, so the posterior is $0.3, 0.3, 0.1, 0.1, 0.1, 0.1$, and the mean is 2.7.

A pass can come only from a looker, and a looker passes only on 3 or 4. So the value is one of those two, with mean 3.5. Note that this leaves the mid where it was, while the range has shrunk to a single unit.

As a check, use the law of total expectation\index{law of total expectation}. A sale and a purchase each have probability $\tfrac{5}{12}$ and a pass $\tfrac16$, and the purchase value is 4.3 by symmetry. Then $\tfrac{5}{12} \cdot 2.7 + \tfrac{5}{12} \cdot 4.3 + \tfrac16 \cdot 3.5 = 3.5$, the original value.

\iqwrong{Assuming the seller surely looked, which gives 1.5. The seller may well be uninformed, and that is what stops the update from going so far.}
\end{solution}

\begin{solution}{iq:iv2:making-a-market:15}
Ask about 100.23 for one lot and 106.40 for 100 lots. A single ask for both sizes is about 106.0.

The value is 110 if the buyer knows, and it averages 100 if the buyer does not. So the ask is $100 + 10 \times$ the probability that the buyer knows, given the size. Size is the signal, and Bayes gives
\[
\Pr(\text{knows} \mid 100) = \frac{0.1 \cdot 0.8}{0.1 \cdot 0.8 + 0.9 \cdot 0.05} = 0.64, \qquad
\Pr(\text{knows} \mid 1) = \frac{0.1 \cdot 0.2}{0.1 \cdot 0.2 + 0.9 \cdot 0.95} \approx 0.023.
\]

A single ask must break even across lots traded, not across traders. Knowing traders send $0.1(0.2 + 80) = 8.02$ lots per arrival and the others send $0.9(0.95 + 5) = 5.355$, so knowing traders send 60\% of the lots, and the ask is about 106.0.

\iqwrong{Averaging across traders. Most traders send one lot and are not informed, but most of the lots are. Those are the lots that cost you money.}
\iqfollow{What does the single ask do to small customers?}{They pay 106 for something worth little over 100, so they leave. And a knowing trader who is priced by size will split into small orders, which is why the signal weakens over time.}
\iqwhy{Showing less size near the mid and more further away follows from this calculation.}
\end{solution}

\begin{solution}{iq:iv2:making-a-market:16}
The lowest ask is 99.21 and the highest bid is 99.19, so the centre is 99.20, and you show 98.20 bid, 100.20 asked. Short 40, you show 99.80 bid, 101.80 asked. The centre moves by 0.02 against each trade.

Selling one lot at $a$ from long 40 changes your score by $a - 100 + 0.01(40^2 - 39^2) = a - 100 + 0.79$, so you accept any $a \ge 99.21$. Buying one lot at $b$ changes your score by $100 - b - 0.01(41^2 - 40^2) = 100 - b - 0.81$, so you accept any $b \le 99.19$.

With position $q$ the centre is $100 - 2 \times 0.01 \times q$. The prices at which you are indifferent are that centre plus or minus 0.01. A purchase takes $q$ up by 1, and the centre down by 0.02. This shifting of both sides together is an inventory skew\index{inventory skew}.

\iqwrong{Lowering only the bid. You want to reduce the position in both ways: to be less likely to buy more, and more likely to sell. Moving both sides keeps your width and shows the same skew.}
\iqfollow{Why would a score have a squared position?}{It is the cost of carrying risk. A position of $q$ on a price with variance $\sigma^2$ has a variance of $q^2\sigma^2$, and a risk-averse holder charges in proportion to that.}
\iqwhy{Many interviewers ask this straight after a trade: ``you are now long, quote me again.'' A candidate who moves the fair value is confusing what they own with what the thing is worth.}
\end{solution}

\begin{solution}{iq:iv2:making-a-market:17}
You expect to earn 9.125 over 10 rounds, against 10 with no limit. In the long run you earn $6/7 \approx 0.86$ per round. With a limit of $\pm L$, it is $2L/(2L + 1)$ per round.

Each trade earns 1, so the question is how many customers you turn away. A position at a limit turns away half of them. Starting from flat, you can first turn someone away in round 4, and stepping the position's probabilities forward exactly gives $73/8$ trades in 10 rounds.

For the long run, the position is a random walk that stays put when it would cross a limit. Every row and every column of its transition matrix sums to 1, so its stationary distribution\index{stationary distribution} is uniform on the $2L + 1$ positions. So you are at a limit with probability $2/(2L + 1)$, and you turn away half the customers then, a share of $1/(2L + 1)$. For $L = 3$, that is one customer in seven.

\iqwrong{Assuming the limit does not bite, as ``we are only at the limit 2 times in 7''. It costs one seventh of your income, and this is with flow that is entirely two-sided.}
\iqfollow{How would you do better?}{Skew before you reach a limit, so that flow brings you back to flat. Here the customers ignore the price, so skew does nothing. Real customers do not.}
\end{solution}

\begin{solution}{iq:iv2:making-a-market:18}
They can lock in a profit if your bid is above 10.8 or your ask is below 10.2. If you hedge the sale at 10.7, you lock in a loss of 0.1. If you do not hedge, you earn 0.2 on average, with a standard deviation of about 2.96.

Three single-die contracts together pay exactly the same as $S$. So they cost 10.8 to buy and fetch 10.2 to sell. If you bid above 10.8, for example 10.9, the counterparty buys the three pieces at 10.8 and sells to you at 10.9, making 0.1 for sure. If you ask below 10.2, they do the reverse. Anywhere from 10.2 to 10.8 is safe from them.

A customer lifts your 10.7, and fair value is 10.5, so you have 0.2 of edge. Buying the pieces at 10.8 gives that up and costs 0.1 more. So a quote inside the pieces' own bid and ask can make money only by carrying the risk for a while: you hold until a customer on the other side arrives, or you decide the risk is worth 0.2.

\iqwrong{Assuming that a hedge is always the right thing. A hedge has a price. Compare what it costs with the risk it removes.}
\iqfollow{Where would you quote if you must hedge every trade at once?}{At or outside 10.2 and 10.8. Then a trade locks in zero or better.}
\end{solution}

\begin{solution}{iq:iv2:making-a-market:19}
$A$ is now worth $28/3 \approx 9.33$ and $B$ is worth $49/6 \approx 8.17$. The best hedge is only $4/23 \approx 0.17$ of a lot of $B$, and it removes about 8.7\% of the variance.

They buy when $X + Y$ is at least 8, which has probability $15/36$. Then $\mathbb{E}[X + Y] = 140/15 = 9.33$. By symmetry $X$ takes half of the excess over 7, so $B$ rises by half as much as $A$: it goes from 7 to 8.17. The cross-update\index{hedge ratio} is the slope of $B$ on $A$, which is $\tfrac12$.

For the hedge, compute given what you now know. You are short $A$, and given $X + Y \ge 8$, the variance of $A$ is $14/9 \approx 1.56$, the variance of $B$ is $161/36 \approx 4.47$, and their covariance is $7/9$. The variance-minimising amount of $B$ is $(7/9)/(161/36) = 4/23$. It removes $(7/9)^2/\left(\tfrac{14}{9} \cdot \tfrac{161}{36}\right) \approx 8.7\%$.

\iqwrong{Using the unconditional ratio of $\tfrac12$. Before any trade that is right, and it removes 25\% of $A$'s variance. Now $A$ is already pinned down to a narrow range, and $B$ still carries all of the noise in $Z$. A hedge of $\tfrac12$ here raises the variance from 1.56 to 1.90.}
\iqfollow{What if they had known only $Y$?}{Then $B$ does not move at all. How far one quote should move when another trades depends on what the trader knows, and you have to guess at that.}
\end{solution}

\begin{solution}{iq:iv2:making-a-market:20}
$h = 1.6/10.8 \approx 0.148$. With a 50\% chance of being hit, $h \approx 0.095$. With 5 customers, $h \approx 0.276$.

Over one cycle you earn $10h$ from customers. If the fast trader gets you, you sell at $f + h$ something now worth $f + 2$, or buy in the mirror image, and either way you lose $2 - h$. So break-even requires
\[
n h - p (M - h) = 0, \qquad h = \frac{p M}{n + p},
\]
with $n = 10$ customers, $p = 0.8$ and a move of $M = 2$. Each customer ends up paying for your share of the news: about $pM/n$ when $n$ is large.

\iqwrong{Comparing $h$ only with the move of 2, and concluding that you need a width of 2. You are not hit on every customer, only once per piece of news. Width is that one hit spread across all the customers.}
\iqfollow{What does this say about speed?}{It is worth money, and you can price it. Cutting $p$ from 0.8 to 0.5 lowers the break-even half-width by 0.053. So you can quote tighter, take more customers, or keep the difference. It is also why a slower dealer widens, or quotes less, before scheduled news.}
\iqwhy{This is the plainest model of being picked off when your quote is stale, a cost that any quoting desk has to cover.}
\end{solution}

\begin{solution}{iq:iv2:making-a-market:21}
Joining earns 0.40 per trader. Posting inside at half-width $h$ earns $h - 0.2$, so $h = 0.5$ earns 0.30, and $h = 0.9$ earns 0.70. If the rival stays put, 0.9 is best. If the rival also moves, quotes fall to a half-width of 0.3 to 0.4, just above the break-even half-width of 0.2.

Per trade at half-width $h$, an uninformed trader pays you $h$, and a knowing trader costs you $2 - h$. Together, $0.9h - 0.1(2 - h) = h - 0.2$. At 99 and 101 that is 0.8 per trade, and sharing gives 0.40. Posting inside beats joining when $h - 0.2 > 0.4$, which is $h > 0.6$, so small improvements win.

When both move, one step takes the whole flow rather than half. At a shared half-width $h$, you earn $(h - 0.2)/2$, and undercutting by 0.1 earns $h - 0.3$. Undercutting pays when $h > 0.4$. It ties at 0.4, and at 0.3 it loses. So the race stops there, one or two steps above break-even, because of the step size. At break-even, width is $2 \times 0.2 = 0.4$, which is the knowing share times the range, $0.1 \times 4$.

\iqfollow{What if a rule lets the first in the queue take the whole trade?}{Then joining is worth less, or nothing if you join behind. Being early in the queue has a value of its own, as it does on a real exchange with price-time priority.}
\end{solution}

\begin{solution}{iq:iv2:making-a-market:22}
$h = 0.8$, a width of 1.6. With one quarter, $h = 0.375$. With three quarters, $h = 4/3 \approx 1.33$.

Someone who has seen a die $d$ values the sum at $d + 7$, so 8 to 13, or $\pm 0.5, \pm 1.5, \pm 2.5$ around 10.5. They trade only when that is outside your quote, and on average they then take the amount by which it exceeds $h$. For $h$ between 0.5 and 1.5, those amounts are 0, $1.5 - h$ and $2.5 - h$, each twice out of six, so the average loss to them is $(4 - 2h)/3$.

The coin toss pays you $h$ on average, so breaking even means
\[
(1 - \alpha)\, h = \alpha \cdot \frac{4 - 2h}{3}.
\]
With $\alpha = \tfrac12$, this gives $h = 0.8$, and that lies in the range assumed. The other two cases solve the same way, each with its own range: $\alpha = \tfrac14$ gives $h$ below 0.5, where the loss is $(4.5 - 3h)/3$, and $\alpha = \tfrac34$ again uses the range from 0.5 to 1.5.

\iqwrong{Using the sum's standard deviation of 2.96. A single seen die tells them much less than that, and their value ranges only $\pm 2.5$, so the width needed is much smaller.}
\end{solution}

\begin{solution}{iq:iv2:making-a-market:23}
The quotes are below. The width first falls below 5 after four sales. It narrows as you learn, although you might expect it to grow as you take more losses. A run of four sales has probability $1/256$, about 0.4\%, if the value is really 100.
\begin{center}\small
\begin{tabular}{@{}lccccc@{}}
Sales so far & 0 & 1 & 2 & 3 & 4 \\ \hline
Bid & 25.00 & 10.00 & 3.57 & 1.22 & 0.41 \\
Ask & 75.00 & 50.00 & 25.00 & 10.00 & 3.57 \\
Width & 50.00 & 40.00 & 21.43 & 8.78 & 3.16
\end{tabular}
\end{center}

A sale has probability $0.75$ if the value is 0 and $0.25$ if it is 100, so each sale divides the odds of 100 by 3. After $n$ sales the odds are $3^{-n}$. The ask is the value given one more purchase, at odds $3^{1-n}$, and the bid is the value given one more sale, at odds $3^{-n-1}$. So the ask after $n$ sales equals the bid after $n - 2$.

Width reflects how much a single trade can still tell you. At the start, the value is entirely unknown. After four sales, you are nearly sure, so a trade has little left to reveal.

\iqwrong{Thinking width grows when flow is one-sided. Your bid falls fast, which is painful, but the width is a separate thing. What widens a real market is new uncertainty, and that is not in this model.}
\iqfollow{And if the value is 0?}{Then the same run has probability $(3/4)^4 \approx 32\%$, so it is an ordinary thing to see. That ratio is why you believe it.}
\iqwhy{After a sharp fall with steady selling, a candidate who widens automatically has confused pain with uncertainty.}
\end{solution}

\begin{solution}{iq:iv2:making-a-market:24}
On a scale of 0 to 1, the ask is $a = \dfrac{(1 + \alpha) - \sqrt{1 - \alpha^2}}{2\alpha}$, and the bid is $1 - a$. At $\alpha = 0.5$, the quote is 36.6 bid, 63.4 asked. As $\alpha$ approaches 1, the ask rises to 100 and the bid falls to 0, so no trade happens: the market breaks down\index{market breakdown}.

The ask equals the expected value given a purchase. A purchase comes from a knowing trader when $V > a$, and from the coin toss at any $V$. So
\[
a = \frac{\alpha \int_a^1 v\,dv + \frac{1 - \alpha}{2} \int_0^1 v\,dv}{\alpha (1 - a) + \frac{1 - \alpha}{2}}
= \frac{\frac{\alpha}{2}(1 - a^2) + \frac{1 - \alpha}{4}}{\alpha(1 - a) + \frac{1 - \alpha}{2}}.
\]
Clearing the fractions gives $\alpha a^2 - (1 + \alpha)a + \frac{1 + \alpha}{2} = 0$, and the smaller root is the one in $[0, 1]$.

Small $\alpha$ gives $a \approx \frac12 + \frac{\alpha}{4}$, so the ask is about $50 + 25\alpha$. For example, $\alpha = 0.1$ gives 52.5, and $\alpha = 0.99$ gives 93.4.

\iqfollow{Why does the market break down only at $\alpha = 1$ here?}{Because any uninformed flow at all pays for a quote wide enough. With a continuum of values, a knowing trader's profit shrinks as the ask nears 100, while the flow keeps its fee. With real costs, it can break down sooner.}
\iqwhy{Liquidity vanishes when informed trading is feared, as it does around a large unknown announcement. This is the simplest model that shows it.}
\end{solution}

\begin{solution}{iq:iv2:making-a-market:25}
The best half-width is $h = 250/21 \approx 11.9$, and it earns about 0.95 per trader. You only break even at $h = 50/7 \approx 7.14$. Above a knowing share of 25\%, no quote is profitable.

A knowing trader takes $(V - 50 - h)^+$ on the upside, and the same on the downside, so with $V$ uniform that averages $(50 - h)^2/100$. An uninformed trader trades with probability $1 - h/20$ and pays $h$ on average. With share $\alpha$,
\[
\pi(h) = (1 - \alpha)\, h \left(1 - \frac{h}{20}\right) - \alpha\, \frac{(50 - h)^2}{100}.
\]
Setting the derivative to zero gives $h = 1\big/\left(\frac{1 - \alpha}{10} + \frac{\alpha}{50}\right)$. At $\alpha = 0.2$, this is $1/0.084 \approx 11.9$. Setting $\pi(h) = 0$ gives $h = 50/7$ for the smaller root.

So the two quotes tell different stories. A market maker with no competition takes 11.9. Competition pushes the half-width down to 7.14, where profit is zero. At $\alpha = \tfrac14$, the best half-width is 12.5, and it earns exactly zero, so at any larger share you should not quote at all.

\iqwrong{Assuming wider is always better. It is better against knowing traders, but you lose customers, and that sets the limit.}
\iqfollow{Why does the market not break down gradually?}{Because the uninformed flow's fee has a ceiling while the knowing traders' take grows. When the most you can make crosses zero, you stop, and you stop all at once.}
\end{solution}

\begin{solution}{iq:iv2:making-a-market:26}
Flat, you show 99.49 bid, 100.51 asked. Long 40, you show 98.69 bid, 100.00 asked. The ask reaches 100 from a position of 26.

The customer's side is a coin toss, so choose each price separately. If a trade changes your penalty by $c$, you maximise $(\delta - c)\,e^{-2\delta}$, whose derivative vanishes at $\delta = \tfrac12 + c$.

Flat, either trade costs 0.01 in penalty, so $\delta = 0.51$. Long 40, a purchase costs $0.01(41^2 - 40^2) = 0.81$, so the bid depth is $1.31$. A sale saves $0.79$, so $c = -0.79$, and the formula says $\delta = -0.29$. But an ask through 100 gets no more fills than an ask at 100, and only lowers the price. So you stop at $\delta = 0$, and earn 0.79. This is the corner of the problem.

Your ask reaches 100 when $0.01(2q - 1) \ge \tfrac12$, so at $q = 26$. Your expected score from the trade is about 0.41, against about 0.18 for the flat quote at the same position.

\iqwrong{Printing the negative depth: quoting below fair value is giving money away, once the probability of a fill has reached one.}
\iqfollow{What is the $\tfrac12$?}{With no inventory cost, you maximise $\delta e^{-k\delta}$, which gives $1/k$. It is the markup that a seller facing that demand curve would charge.}
\end{solution}

\begin{solution}{iq:iv2:making-a-market:27}
With no signal, quote 97 bid, 103 asked. With the signal saying 110, quote about 103.66 bid, 107.63 asked, a width of 3.97 against 6. A signal correct only 60\% of the time leaves about 98.94 to 104.72, a width of 5.78.

A public signal changes your prior, and the same Bayes step then gives the quote. Start with the signal, so the chance of 110 is 0.8, and treat a buy as a further observation with likelihoods $0.65$ and $0.35$. Then
\[
\Pr(110 \mid \text{buy}) = \frac{0.8 \cdot 0.65}{0.8 \cdot 0.65 + 0.2 \cdot 0.35} \approx 0.881,
\]
which gives an ask of $90 + 20 \times 0.881 = 107.63$. A sale reverses the likelihoods, and that gives 103.66.

Note that the quote is not centred on the new fair value, which is $0.8 \times 110 + 0.2 \times 90 = 106$. The mid is at 105.64, because a trade can move you further towards 90 than towards 110.

\iqwrong{Thinking the share of knowing traders has fallen. It has not. But public knowledge shrinks what they know beyond you, so each of their trades takes less.}
\iqfollow{What would a perfect signal do?}{It would take the width to zero, as no one knows anything you do not. The value of a signal to a market maker is lower adverse selection, which means tighter quotes and more flow.}
\end{solution}

\begin{solution}{iq:iv2:making-a-market:28}
At $m = 0.5$ you lose about 0.236 per deal. Break-even is $m = 35/36 \approx 0.97$. With three dealers, it is $35/24 \approx 1.46$.

You sell only when your own die is the lower one. Your ask is $d_1 + 3.5 + m$, and the contract pays $d_1 + d_2$, so your profit on a sale is $3.5 + m - d_2$, and winning selects a high $d_2$. If your die is lower you win everything, and if they tie you win half. So given $d_2$, your share is $(d_2 - \tfrac12)/6$, and
\[
\text{profit} = \frac{1}{36} \sum_{d_2 = 1}^6 \left(3.5 + m - d_2\right) \left(d_2 - \tfrac12\right) = \frac{18m - 17.5}{36}.
\]
This vanishes at $m = 35/36$. Among the deals where your die is strictly lower, the rival's die averages $14/3$, not 3.5. This is the winner's curse\index{winner's curse}: you win the trade when you have underestimated.

With three dealers, the same count over 216 outcomes gives $m = 35/24$. More rivals make it worse, because winning says more.

\iqwrong{Using the sum's unconditional value and adding a small edge. That edge is entirely eaten by the fact that you only sell when the other die is high.}
\iqfollow{Is this the same as adverse selection from a customer?}{It has the same structure, but the customer here knows nothing. The information comes from your rival, through the customer's choice between you.}
\end{solution}

\begin{solution}{iq:iv2:making-a-market:29}
Quote 25 bid, 75 asked in round 1. After a purchase, quote 50 bid, 83.33 asked. After a purchase and then a sale, quote 50 bid, 50 asked. With a new person each round, the same three quotes are 25 to 75, 50 to 90, and 25 to 75.

There are four states: knows or does not, and 0 or 100, each of probability $\tfrac14$. A purchase has likelihood 1 from a knower at 100, 0 from a knower at 0, and $\tfrac12$ otherwise. So after one purchase the states have probabilities: knower at 100, $\tfrac12$; no knowledge at 100, $\tfrac14$; no knowledge at 0, $\tfrac14$. The value is 75.

The next ask is the value given a second purchase, which reweights those to $\tfrac12 : \tfrac18 : \tfrac18$, giving $100 \times \tfrac{5/8}{6/8} = 83.33$. The next bid is the value given a sale. A knower who has just bought cannot sell, so a sale proves they know nothing, and the value goes back to exactly 50. Since a person who reverses is uninformed and their trades now say nothing, you need no width at all.

\iqwrong{Treating one person as a new draw every round. A second trade by the same person is much less informative than a trade by someone new, so the ask is lower than 90, and a reversal is decisive.}
\iqfollow{What happens to the bid if they keep buying?}{It stays at 50. The ask rises from 83.33 to 90, then towards 100, since each purchase makes it likelier they know.}
\end{solution}

\begin{solution}{iq:iv2:making-a-market:30}
The unconditional value is $33/4 = 8.25$. The best quote is 7.5 bid, 9.5 asked, which loses 0.05 on average. It is centred on 8.5, above the value. If they buy, the contract is worth 10, and if they pass, it is worth $145/18 \approx 8.06$. They never sell.

Their value depends only on the card they saw:
\begin{center}\small
\begin{tabular}{@{}l*{10}{c}@{}}
Card & 1 & 2 & 3 & 4 & 5 & 6 & 7 & 8 & 9 & 10 \\ \hline
Value & 7.67 & 7.67 & 7.67 & 7.69 & 7.78 & 7.94 & 8.22 & 8.64 & 9.22 & 10.00
\end{tabular}
\end{center}
Low cards barely matter, because one of the other two is usually higher. The values are bunched at the bottom and spread at the top.

A quote centred on 8.25, at 7.25 to 9.25, has no seller either, but sells to card 10 at 9.25, losing $0.75/10 = 0.075$. Moving up to 7.5 to 9.5 cuts that to 0.05. At 7.75 to 9.75, cards 1 to 4 begin to sell to you, and the loss returns to about 0.056. With a skewed distribution of their value, the best quote is not centred on the mean. As a check, $0.9 \times 145/18 + 0.1 \times 10 = 8.25$.

\iqwrong{Centring on the value by habit. A quote protects you against the other side's value, so it should sit around the distribution of that value, which is not symmetric here.}
\end{solution}

\begin{solution}{iq:iv2:making-a-market:31}
If they bought $S$, then $S$ is worth 9, and $D$ stays at 7. If they buy both, both are worth 9, and they saw $P$. If they buy $S$ and sell $D$, then $S$ is worth 9 and $D$ is worth 5, and they saw $Q$. You lose 2 on average per deal.

Seeing $P = p$, they value both contracts at $p + 3.5$. Seeing $Q = q$, they value $S$ at $q + 3.5$ and $D$ at $10.5 - q$. Either way they buy $S$ when the seen die is 5 or 6. So a purchase of $S$ raises its value by $5.5 - 3.5 = 2$, whichever die it was. For $D$, it is up 2 or down 2 with equal chance, so unchanged on average.

Their trade on $D$ then tells you which die. Buying both means $P$ was the seen die, so $D$ rises to 9. Buying $S$ and selling $D$ means it was $Q$, so $D$ falls to 5.

If the seen die is 3 or 4, they pass on both, which happens with probability $\tfrac13$. Otherwise you lose 1.5 on each of two trades, and $\tfrac23 \times 3 = 2$.

\iqwrong{Moving $D$ with $S$ because they are related. Their covariance is zero, as they have equal variance from $P$ and $Q$ with opposite signs on $Q$, so one trade in $S$ alone tells you nothing about $D$ on average.}
\iqfollow{So how should you update on $D$ after seeing only the $S$ trade?}{Keep the mid, but widen: $D$ will now be 5 or 9, and you do not yet know which. Your uncertainty about it has grown though your best guess has not moved.}
\end{solution}

\begin{solution}{iq:iv2:making-a-market:32}
Quote through the middle each time, which is a binary search\index{binary search}: start at 5 bid, 6 asked, then split what remains. The expected total loss is $24/10 = 2.4$. With numbers 1 to 16, it is $4.75$. Walking up from 1 bid, 2 asked costs 7.

Start at 5 and 6. If they pass, the number is 5 or 6, and you never lose again, so you show 5 bid, 6 asked forever. If they sell, the number is from 1 to 4, and you have lost $5 - x$. Next quote 2 bid, 3 asked. If the number is 1, they sell again and you lose 1, and then you know it. If it is 4, they buy, and you lose 1. So the total losses for the numbers 1 to 10 are:
\begin{center}\small
\begin{tabular}{@{}l*{10}{c}@{}}
Number & 1 & 2 & 3 & 4 & 5 & 6 & 7 & 8 & 9 & 10 \\ \hline
Total loss & 5 & 3 & 2 & 2 & 0 & 0 & 2 & 2 & 3 & 5
\end{tabular}
\end{center}
That sums to 24. The first quote costs 2.0 on average, and the second only 0.4 more. A dynamic program\index{dynamic programming} over all sets of remaining numbers confirms that no other quoting does better, and that more rounds cost nothing once you have the number.

Walking upwards loses far more. If the number is 10, you lose 8, then 6, 4 and 2, for a total of 20, and the average over all numbers is 7. Every quote is on the wrong side of the same number.

\iqfollow{What would you pay to be told the number?}{Up to 2.4, because that is what the information costs you to find out by quoting.}
\iqwhy{After a trade against you, an aggressive move of the quote costs less in total than a slow walk. A candidate who moves by a little each time pays the most.}
\end{solution}

\begin{solution}{iq:iv2:making-a-market:33}
Take a side down when one more lot would add more penalty than it can repay. The largest position you will go to, by rounds left from 1 to 10, is 2, 3, 3, 4, 4, 5, 5, 5, 6, 6. The best expected score is 8.41, against 8.00 for never taking a side down, and 8.35 for the best fixed limit, which is $\pm 3$.

Work back from the end by backward induction\index{backward induction}. With one round left, a purchase at position $q$ gains $1 - 0.2(2q + 1)$, which is positive for $q = 1$ and exactly zero for $q = 2$. So you buy up to 2 and no further. With two rounds left, compute the value of each position with one round left, and compare again. A purchase now is worth more than it would be at the end, because one more trade may take you back down. So you buy up to 3.

Each step back adds time for flow to return you towards flat, so the allowed position grows with the rounds left.

Never taking a side down earns one per round, or 10, minus $0.2 \times \mathbb{E}[q^2] = 0.2 \times 10$, so 8.00. A fixed limit of $\pm 3$ gets most of the gain but is too tight early and too loose at the end.

\iqwrong{Treating a limit as one fixed number. It should depend on how long you have to work out of the position. A position with a day to unwind is not the same as one at the close.}
\iqfollow{What if the customers respond to price?}{You can then skew your quote instead of refusing trades. Taking a side down is an extreme skew, and a price does the same job more smoothly.}
\end{solution}

\begin{solution}{iq:iv2:making-a-market:34}
$\lambda = \sigma_v/(2\sigma_u) = 0.01$ per share. The insider's order is $x = (\sigma_u/\sigma_v)(V - 100) = 50(V - 100)$, and the expected profit is $\sigma_v\sigma_u/2 = 100$. After the price is set, half the variance of $V$ remains, so 2 out of 4. This is the one-period Kyle model\index{Kyle model}.

The insider chooses $x$ to maximise $x(V - 100 - \lambda x)$, since the noise averages to zero. That gives $x = \beta(V - 100)$ with $\beta = 1/(2\lambda)$. You set the price to the expected value given the flow, so $\lambda$ is the regression slope of $V$ on $y$:
\[
\lambda = \frac{\mathrm{Cov}(V, y)}{\mathrm{Var}(y)} = \frac{\beta \sigma_v^2}{\beta^2 \sigma_v^2 + \sigma_u^2}.
\]
Substituting $\beta = 1/(2\lambda)$ gives $4\lambda^2\sigma_u^2 = \sigma_v^2$, which is the result.

The insider's order has the same variance as the noise, because $\beta\sigma_v = \sigma_u$. It hides in the crowd exactly. So half of the variance of $y$ is signal, and the price moves half of the way to the truth. The insider's expected profit is $\mathbb{E}\left[\beta (V-100)^2\right]/2 = 50 \times 4/2 = 100$, and the noise traders pay for it: each order $u$ moves the price against itself by $\lambda u$, which costs them $\lambda\sigma_u^2 = 100$ on average.

\iqwrong{Having the insider trade until the price reaches the true value. The order moves the price, so the best size is only half of the one that would remove the edge.}
\iqfollow{What happens if noise doubles?}{$\lambda$ halves, the insider's order and profit double, and the price becomes no less informative. More noise pays the insider and does not change the variance left.}
\end{solution}

\begin{solution}{iq:iv2:making-a-market:35}
The bid is about 98.35, and the ask is about 100.05. The width is $\gamma\sigma^2 + \frac{2}{\gamma}\ln\left(1 + \frac{\gamma}{k}\right) \approx 0.40 + 1.29 = 1.69$, and the centre is $100 - \gamma\sigma^2 q = 99.2$, with $\gamma = 0.1$, $\sigma^2 = 4$, $k = 1.5$ and $q = 2$. As $\gamma$ goes to 0, each side goes to $1/k$ from 100, so about 99.33 to 100.67.

With no fill, exponential utility of a normal gives $U_0 = -\exp\left(-\gamma(x + 100q) + \tfrac12\gamma^2\sigma^2 q^2\right)$, where $x$ is your cash. If your bid at depth $\delta$ is taken, you pay $100 - \delta$ and hold $q + 1$, so
\[
U = U_0 \exp\left(-\gamma(\delta - d_b)\right), \qquad d_b = \tfrac12\gamma\sigma^2(2q + 1).
\]
Here $d_b$ is the inventory cost of that purchase. Expected utility on that side is
\[
U_0\left[1 + e^{-k\delta}\left(e^{-\gamma(\delta - d_b)} - 1\right)\right].
\]
Since $U_0$ is negative, you want the bracket small. Setting $u = \delta - d_b$, the first-order condition gives $e^{-\gamma u} = k/(k + \gamma)$, which is
\[
\delta = d_b + \frac{1}{\gamma}\ln\left(1 + \frac{\gamma}{k}\right).
\]
For the ask, the same steps give $d_a = \tfrac12\gamma\sigma^2(1 - 2q)$. So $\delta_b = 1.0 + 0.645$ and $\delta_a = -0.6 + 0.645 = 0.045$. Holding 2 shares, you barely charge to sell, and you charge a lot to buy.

This is the one-period form of the Avellaneda-Stoikov quote\index{Avellaneda-Stoikov model}. The shift depends on your position and the risk, and the width on the risk and the customers' sensitivity to price.

\iqwrong{Adding the inventory term to one side only. It moves both sides, and $d_b + d_a = \gamma\sigma^2$ is the part it adds to the width.}
\end{solution}
